\chapter{The Dual-License Engine: BUSL-1.1 Plus Commercial}
\label{ch:licensing}

\section{The structure}

wasm4pm is licensed under the Business Source License 1.1 with a commercial
license available for production use that falls outside the BUSL grant. After
the Change Date the code converts automatically to AGPL-3.0. This is the
MariaDB/HashiCorp/Sentry pattern, and it creates a three-segment funnel:

\begin{description}[leftmargin=2em]
  \item[Free tier (BUSL-compliant use).] Evaluation, development, internal
    tooling, academic research. This population is the adoption engine ---
    it produces GitHub stars, npm downloads, bug reports, and the analysts
    who later champion a paid deployment.
  \item[Commercial tier.] Anyone embedding wasm4pm in a product or service
    offered to third parties, or whose production use exceeds the BUSL
    additional-use grant. This is the revenue engine.
  \item[Post-Change-Date tier.] Code older than the change window is AGPL ---
    a strong-copyleft poison pill for proprietary embedders that keeps the
    commercial license attractive even for old versions.
\end{description}

\section{Why BUSL fits this product specifically}

BUSL is usually defended as cloud-provider insurance. For wasm4pm there is a
sharper argument: \textbf{the product is a trust primitive, and trust
primitives must be source-available.} A receipt chain whose verifier is
closed-source is marketing, not evidence. A regulated customer's auditors
must be able to read the BLAKE3 hashing path and the conformance scoring
code. Fully proprietary licensing would destroy the core value proposition;
fully permissive licensing (the previous MIT/Apache posture) would invite a
hyperscaler to operate ``Process Mining as a Service'' on the project's own
code. BUSL is the narrow corridor between those failure modes.

\section{Pricing axes}

The architecture suggests charging on axes that scale with value received,
not with seats:

\begin{center}
\begin{tabularx}{\textwidth}{lX}
\toprule
\textbf{Axis} & \textbf{Rationale} \\
\midrule
Receipts retained / verified per month & Directly proportional to audit value; metered by the Supabase sync layer that already counts them. \\
Deployment profiles unlocked & Browser-only free; edge/IoT/fog profiles commercial --- the size-optimized builds are real engineering, easily gated. \\
Algorithm tiers & Core discovery free; conformance, OCEL 2.0 querying, and the cognition layer commercial. \\
Embedding rights & OEM license for ISVs shipping wasm4pm inside their own products --- the classic dual-license cash cow. \\
\bottomrule
\end{tabularx}
\end{center}

\section{The CalVer signal}

Versioning by calendar date (\texttt{v26.6.9} = June 9, 2026) is itself a
licensing instrument: the BUSL Change Date is defined per release, so CalVer
makes the conversion schedule legible at a glance. A customer can see
precisely when any given version becomes AGPL, which converts license
anxiety into a predictable procurement calculation.
